\documentclass[conference]{IEEEtran}
% Requires IEEEtran.cls (obtain from IEEE/CTAN; provided by TeX Live, MiKTeX, Overleaf).
\IEEEoverridecommandlockouts
\usepackage{cite}
\usepackage{graphicx}
\usepackage{booktabs}
\usepackage{url}
\usepackage{hyperref}

\begin{document}

\title{Cybersecurity Posture of Sri Lankan Government Digital Infrastructure: A
Ten-Institution Public-Evidence-Based Study}

\author{\IEEEauthorblockN{Vinod Ravindu Perera}
\IEEEauthorblockA{Independent Researcher\\
Sri Lanka\\
Repository: \url{https://github.com/Perera1325/sri-lanka-gov-security-research}}}

\maketitle

\begin{abstract}
Sri Lankan government institutions have experienced a series of publicly reported website and
infrastructure security incidents, yet little is published about how these events are followed
up or about the state of the institutions' public web security. This paper presents an
exploratory, public-evidence-based study of ten Sri Lankan government digital institutions.
Using only publicly available information and passive browser access, it compiles and
classifies a dataset of 20 historical incident records; maps the 32 requirements of the Sri
Lanka CERT web-application security guideline against external verifiability; and reviews the
academic literature on AI-generated-code security mapped to the OWASP Top 10, without claiming
any studied website uses such code. Six website compromises affecting five institutions are
confirmed from public reporting; for none of the thirteen confirmed incidents was a published
investigation outcome found. Only seven of the 32 requirements are informable from public
evidence, and twenty-one require authorization or internal access. No active testing was
performed and no current exploitable weakness is claimed. The contribution is a transparent,
reproducible method and a baseline dataset for public-accountability research.
\end{abstract}

\begin{IEEEkeywords}
government cybersecurity, public-sector web security, Sri Lanka, security governance, incident
analysis, OWASP, AI-generated code, responsible research
\end{IEEEkeywords}

\section{Introduction}
Government institutions increasingly deliver services through web applications, and the security
of that infrastructure is a public concern. In Sri Lanka, several government websites have been
reported compromised, defaced or disrupted, and Sri Lanka CERT has issued security guidelines
for government organizations. Independent assessment of government systems normally requires
written authorization, because unauthorized testing is unlawful under the Computer Crimes Act
No.~24 of 2007~\cite{slcca2007}. This leaves a gap between visible incidents and verifiable
governance information. We address it with a constrained, public-evidence method that classifies
every result by a fixed evidence scheme and is explicit about the boundary between what public
evidence can and cannot establish. We contribute (i) a transparent, classified incident dataset,
(ii) a verifiability mapping of the national web-application guideline, and (iii) an OWASP-mapped
review of AI-generated-code risk as procurement context, with no link asserted between any site
and AI-generated code.

\section{Background and Related Work}
Sri Lanka CERT publishes website and web-application security guidelines for government
organizations~\cite{slcert2022website,slcert2022webapp}, providing a public yardstick; the
Computer Crimes Act~\cite{slcca2007} constrains research to passive, public-evidence methods
absent authorization. Where authorized testing is unavailable, open-source intelligence observes
configuration and public record, not exploitability---a limitation treated here as a design
constraint. A growing peer-reviewed literature examines AI-generated-code security, reporting
frequent weakness patterns, developer over-confidence, memorized secrets, and package
hallucination~\cite{pearce2022asleep,perry2023do,sandoval2023lost,fu2025security,vero2025baxbench,spracklen2025package,huang2024yourcode,klemmer2024using}.
These concern software in general and are used here only as policy context.

\section{Research Questions}
\textbf{RQ1:} What cybersecurity incidents have been publicly documented across the selected
institutions, and how should they be classified? \textbf{RQ2:} How do documented
security-governance requirements compare with the evidence available about these institutions
from the outside? \textbf{RQ3:} What AI-assisted-development risks are relevant to government
digital infrastructure as policy context?

\section{Methodology}
The method is passive and public-evidence-based (Fig.~\ref{fig:method}). We define scope (ten
\texttt{gov.lk} institutions, public pages only; no scanning, fuzzing, injection or login
attempts); verify each official domain against independent government sources; collect public
records; classify every result by a seven-label evidence scheme (Table~\ref{tab:evid}); map the
CERT guideline by verifiability; and review AI-code literature mapped to OWASP, kept separate
from the institution findings. Every count is regenerated from the data files, and structural and
safety invariants are enforced by a validation script.

\begin{figure}[t]\centering
\includegraphics[width=0.9\columnwidth]{methodology-flowchart.pdf}
\caption{Research methodology. Render from \texttt{figures/methodology-flowchart.md}.}
\label{fig:method}\end{figure}

\section{Institution Selection}
Ten institutions span public security, administration, records, health regulation, education,
taxation, immigration, transport and digital-economy policy (Table~\ref{tab:inst}). None is yet
fully verified: seven are PRELIMINARY and three PARTIAL (Fig.~\ref{fig:verif}). Observation would
not begin until verification is complete.

\section{Historical Cyber Incident Analysis (RQ1)}
The dataset has 20 records: 13 confirmed historical incidents, 2 unverified reports, 5 with no
public evidence (Table~\ref{tab:incident}, Fig.~\ref{fig:incident}). By system affected, ten
concern a website, four a backend system, three are impersonation, two are
shared-infrastructure events, and one concerns social-media accounts. Six website compromises
are confirmed, affecting five of the ten institutions; one is an explicit defacement and the
rest are unauthorized access or alteration. Impersonation and shared-infrastructure events are
reported separately and not counted as website vulnerabilities. Critically, for none of the
thirteen confirmed incidents was a published investigation outcome found: incidents are
reported and inquiries announced, but outcomes do not appear in the public record.

\section{Sri Lanka CERT Governance Analysis (RQ2)}
The CERT web-application guideline sets 32 requirements. Only seven are informable, even partly,
by passive observation; four more only via an official statement; five require written
authorization; sixteen are not publicly verifiable (Table~\ref{tab:cert},
Fig.~\ref{fig:cert}). Twenty-one of 32 cannot be checked without authorization. The one fully
verifiable requirement---a \texttt{gov.lk} domain---is met by all ten. Most of the national
standard is thus externally invisible, which bounds what any public-evidence study can conclude.

\section{AI-Assisted Development Security Risks (RQ3)}
As procurement context, we map eight verified papers' reported weakness classes to the OWASP Top
10 (Table~\ref{tab:owasp}). The literature reports frequent weaknesses in generated code,
over-confidence, memorized credentials, and package hallucination; results on finished-code
security are mixed, and human review that does not trust the output is the key control. Three
OWASP categories (A05, A09, A10) are uncovered by the sources; this is stated rather than filled.
No source links any studied website to AI-generated code, and site appearance is never treated
as evidence of how it was built.

\section{Discussion}
Three themes emerge: incidents are real but published accountability is thin; most of the
security standard is externally invisible; and the need for authorization to verify most
requirements is itself an argument for authorized, coordinated assessment rather than
unauthorized testing or silence.

\section{Recommendations}
Publish incident outcomes; publish basic non-sensitive security status; provide a lawful
authorized-assessment channel; and require, in procurement, human security review of generated
code, dependency verification and secrets handling.

\section{Limitations and Threats to Validity}
No technical assessment was performed; the evidence register is empty and no current weakness is
claimed. Historical incidents are not current weaknesses. Most records rest on secondary news,
not primary statements. No domain is fully verified. Absence of evidence is not evidence of
absence. The ten institutions are not a representative sample. The AI-code literature concerns
software in general, and three OWASP categories are uncovered.

\section{Conclusion}
From public evidence alone, six website compromises across five of ten institutions are
confirmed, no published investigation outcomes are found, and only a small part of the national
web-application standard is externally verifiable. The contribution is a transparent,
reproducible, legally constrained method and an open dataset, with an honest account of its
limits. The next step is authorized, coordinated assessment.

\section*{Acknowledgment}
Parts of this research infrastructure and draft were prepared with AI assistance under the
author's direction and review. Data-classification decisions, the passive-only restriction, and
reference verification are the author's responsibility.

% Tables (port generated Markdown tables from ../tables/ here; keep numbers identical).
\begin{table}[t]\caption{Selected institutions (see \texttt{tables/table1}).}\label{tab:inst}
\centering\footnotesize\begin{tabular}{@{}ll@{}}\toprule Count & Value\\\midrule
Institutions & 10\\ With a confirmed website incident & 5\\ Fully verified domains & 0\\\bottomrule
\end{tabular}\end{table}

\begin{table}[t]\caption{Incident classification (see \texttt{tables/table2}).}\label{tab:incident}
\centering\footnotesize\begin{tabular}{@{}lr@{}}\toprule System type & Records\\\midrule
Website (6 confirmed) & 10\\ Other/backend & 4\\ Impersonation & 3\\ Shared infrastructure & 2\\
Social media & 1\\ \textbf{Total} & \textbf{20}\\\bottomrule\end{tabular}\end{table}

\begin{table}[t]\caption{CERT guideline verifiability (see \texttt{tables/table4}).}\label{tab:cert}
\centering\footnotesize\begin{tabular}{@{}lr@{}}\toprule Class & Requirements\\\midrule
Passive observation & 7\\ Official statement only & 4\\ Written authorization & 5\\
Not publicly verifiable & 16\\ \textbf{Total} & \textbf{32}\\\bottomrule\end{tabular}\end{table}

\begin{table}[t]\caption{Evidence labels (see \texttt{tables/table6}).}\label{tab:evid}
\centering\footnotesize\begin{tabular}{@{}p{0.9\columnwidth}@{}}\toprule
Seven labels separate confirmed incidents, requirements, observed configuration, confirmed
weakness, unverified risk, not-tested, and no-evidence. Labels for the researcher's own
observations are unused (no site observed).\\\bottomrule\end{tabular}\end{table}

\begin{table}[t]\caption{AI-code to OWASP (see \texttt{tables/table5}).}\label{tab:owasp}
\centering\footnotesize\begin{tabular}{@{}lr@{}}\toprule OWASP 2021 & Covered?\\\midrule
A01--A04, A06--A08 & Yes\\ A05, A09, A10 & Not covered by sources\\\bottomrule\end{tabular}\end{table}

\begin{figure}[t]\centering\includegraphics[width=0.95\columnwidth]{fig1-incident-classification.pdf}
\caption{Incident records by system type (n=20).}\label{fig:incident}\end{figure}
\begin{figure}[t]\centering\includegraphics[width=0.8\columnwidth]{fig2-domain-verification.pdf}
\caption{Domain verification status (n=10).}\label{fig:verif}\end{figure}
\begin{figure}[t]\centering\includegraphics[width=0.95\columnwidth]{fig3-cert-verifiability.pdf}
\caption{External verifiability of 32 CERT requirements.}\label{fig:cert}\end{figure}

\bibliographystyle{IEEEtran}
\bibliography{references}

\end{document}
